% =====================================================================
%  MSCA-PF Part B — formatting enforced automatically (edit rarely)
%  Rules encoded here come from the official Application form (V5.0,
%  27.03.2026), section 3.4 "Page limit and formatting":
%    - A4 page size
%    - all body-text margins >= 15 mm  (set in _quarto.yml geometry)
%    - reference font: Times New Roman / Nimbus Roman No. 9 L
%      -> we use TeX Gyre Termes, the free, metric-compatible equivalent
%    - minimum body font size 11 pt, single line spacing (set in YAML)
%    - tables also >= 11 pt (we never shrink table text)
%    - headers/footnotes/captions may be smaller but >= 8 pt
%    - footer: "Part B - Page X of Y"
%  You normally never need to touch this file. Just write in the
%  sections/ files. Formatting takes care of itself.
% =====================================================================

% ---- Body font: TeX Gyre Termes (= Times New Roman / Nimbus Roman No. 9 L) ----
\usepackage[T1]{fontenc}
\usepackage{tgtermes}          % free, embedded, metric-compatible with Times New Roman

% ---- Let common Unicode symbols "just work" under pdfLaTeX ----
% So you can type >=, <=, x, +-, ->, degree, micro, Greek letters etc. straight
% into the text without build errors. (If you need lots of exotic Unicode,
% switch pdf-engine to lualatex in _quarto.yml and delete this block.)
\usepackage{textcomp}
\usepackage{amssymb}
\usepackage{newunicodechar}
\newunicodechar{≥}{\ensuremath{\geq}}
\newunicodechar{≤}{\ensuremath{\leq}}
\newunicodechar{×}{\ensuremath{\times}}
\newunicodechar{±}{\ensuremath{\pm}}
\newunicodechar{≈}{\ensuremath{\approx}}
\newunicodechar{≠}{\ensuremath{\neq}}
\newunicodechar{→}{\ensuremath{\rightarrow}}
\newunicodechar{←}{\ensuremath{\leftarrow}}
\newunicodechar{↔}{\ensuremath{\leftrightarrow}}
\newunicodechar{⇒}{\ensuremath{\Rightarrow}}
\newunicodechar{·}{\textperiodcentered}
\newunicodechar{°}{\textdegree}
\newunicodechar{µ}{\textmu}
\newunicodechar{μ}{\ensuremath{\mu}}
\newunicodechar{α}{\ensuremath{\alpha}}
\newunicodechar{β}{\ensuremath{\beta}}
\newunicodechar{γ}{\ensuremath{\gamma}}
\newunicodechar{δ}{\ensuremath{\delta}}
\newunicodechar{κ}{\ensuremath{\kappa}}
\newunicodechar{λ}{\ensuremath{\lambda}}
\newunicodechar{σ}{\ensuremath{\sigma}}
\newunicodechar{✓}{\ensuremath{\checkmark}}
\newunicodechar{•}{\textbullet}

% ---- Footer "Part B - Page X of Y" (Y = actual total pages) ----
\usepackage{fancyhdr}
\usepackage{lastpage}
\pagestyle{fancy}
\fancyhf{}
\renewcommand{\headrulewidth}{0pt}
\fancyfoot[C]{\fontsize{9}{11}\selectfont Part B - Page \thepage{} of \pageref{LastPage}}
% Make sure the first page (which LaTeX styles as "plain") also shows the footer
\fancypagestyle{plain}{%
  \fancyhf{}%
  \renewcommand{\headrulewidth}{0pt}%
  \fancyfoot[C]{\fontsize{9}{11}\selectfont Part B - Page \thepage{} of \pageref{LastPage}}%
}
% To print "of 10" (the page limit) instead of the actual total, replace
% \pageref{LastPage} above with the number 10.

% ---- Headings: bold, legible, compact (non-body text, so size may differ) ----
\usepackage{titlesec}
\titleformat{\section}{\normalfont\large\bfseries}{\thesection}{0.5em}{}
\titleformat{\subsection}{\normalfont\normalsize\bfseries}{\thesubsection}{0.5em}{}
\titleformat{\subsubsection}{\normalfont\normalsize\bfseries\itshape}{\thesubsubsection}{0.5em}{}
\titlespacing*{\section}{0pt}{8pt}{4pt}
\titlespacing*{\subsection}{0pt}{6pt}{3pt}
\titlespacing*{\subsubsection}{0pt}{4pt}{2pt}

% ---- Lists: tighter but readable (saves precious page space) ----
\usepackage{enumitem}
\setlist{nosep, topsep=2pt, partopsep=0pt, leftmargin=1.4em}

% ---- Tables: keep body size (never below 11 pt) and give rows a little air ----
\renewcommand{\arraystretch}{1.15}

% ---- Captions a touch smaller (allowed: >= 8 pt), still clearly legible ----
\usepackage{caption}
\captionsetup{font=small, labelfont=bf, skip=4pt}

% ---- Sensible microtype spacing to fit more text without shrinking font ----
\usepackage{microtype}
